% ch3_answers.tex -- Quick check answers for Module 7, Chapter 3.
\begin{chapteranswers}{3}{Light as a Particle}
  % ---- 3.1 Black body radiation and Planck's quanta
  \item Colour is not the point (black paint still reflects some light). A black
        body is an ideal object that absorbs \emph{all} the EM radiation that falls on
        it, at every wavelength, and reflects none (1 mark). It is also a perfect
        emitter: its radiation depends only on its temperature, so a hot black body,
        like a star, can look white (1 mark).
  \item $T = 600 + 273 = \qty{873}{K}$. Then
        $\lambda_\mathrm{max} = \dfrac{b}{T} = \dfrac{\num{2.898e-3}}{873}
        = \qty{3.32e-6}{m} = \qty{3.32}{\micro m}$.
        Answer: \final{$\lambda_\mathrm{max} = \qty{3.3e-6}{m}$}, in the
        \textbf{infrared} (longer than \qty{700}{nm}). The element still glows dull
        red, because the short-wavelength side of its curve just reaches the red end
        of the visible band.
  \item $T = \dfrac{b}{\lambda_\mathrm{max}} = \dfrac{\num{2.898e-3}}{\num{1.05e-6}}
        = \qty{2760}{K}$. In Celsius: $2760 - 273 = \qty{2487}{\degreeCelsius}$.
        Answer: \final{$T = \qty{2.76e3}{K}$} (about \qty{2.49e3}{\degreeCelsius}).
  \item $\lambda = \qty{400}{nm} = \qty{4.00e-7}{m}$ (nm to m).\\[2pt]
        \eqlines{E &= \frac{hc}{\lambda} = \frac{\num{6.626e-34} \times \num{3.00e8}}{\num{4.00e-7}}
        = \qty{4.970e-19}{J}\\
        E &= \frac{\num{4.970e-19}}{\num{1.602e-19}} = \qty{3.102}{eV} \quad\text{(J to eV)}}\\[3pt]
        Answer: \final{$E = \qty{4.97e-19}{J} = \qty{3.10}{eV}$}.
  \item Classical physics agreed with experiment at \textbf{long} wavelengths (the
        infrared side). At short wavelengths it predicted intensity rising without
        limit, so infinite total energy. Real curves fall to zero there. The failure
        is on the ultraviolet side of the curve, so it was called the ultraviolet
        catastrophe.
  % ---- 3.2 Photoelectric effect
  \item $K_\mathrm{max} = qV_0 = \num{1.602e-19} \times 0.85 = \qty{1.362e-19}{J}$.
        In eV the number is the same as the voltage.
        Answer: \final{$K_\mathrm{max} = \qty{0.85}{eV} = \qty{1.4e-19}{J}$}.
  \item The vacuum stops the photoelectrons from hitting air molecules. Collisions
        would slow them down or stop them before the anode, so the current and $V_0$
        would be wrong. The quartz window lets ultraviolet light in; ordinary glass
        absorbs ultraviolet.
  \item (a) The photocurrent increases, because more photoelectrons are released
        each second. (b) The stopping voltage does not change. The wave model
        predicted (a) correctly. It wrongly predicted that $V_0$ would increase,
        because it said brighter light gives each electron more energy.
  \item Any two of these, each with its wave-model prediction.
        \textbf{Threshold frequency:} no electrons leave below $f_0$, however bright
        the light; the wave model predicted that any frequency works if the light is
        bright enough. \textbf{No time delay:} electrons leave at once, even in dim
        light; the wave model predicted a delay while each electron builds up energy.
        \textbf{$K_\mathrm{max}$ depends on frequency, not brightness:} the wave
        model predicted that brighter light gives faster electrons.
  % ---- 3.3 Photon model and graphs
  \item \eqlines{\phi &= 2.28 \times \num{1.602e-19} = \qty{3.653e-19}{J} \quad\text{(eV to J)}\\
        \lambda_0 &= \frac{hc}{\phi} = \frac{\num{6.626e-34} \times \num{3.00e8}}{\num{3.653e-19}}
        = \qty{5.44e-7}{m} = \qty{544}{nm}}\\[3pt]
        Answer: \final{$\lambda_0 = \qty{5.44e-7}{m}$}. Yellow light (\qty{589}{nm})
        does \textbf{not} release electrons: \qty{589}{nm} is longer than
        \qty{544}{nm}. Each of its photons has only this much energy:\\[2pt]
        \eqlines{E &= \frac{hc}{\lambda} = \frac{\num{6.626e-34} \times \num{3.00e8}}{\num{5.89e-7}}
        = \qty{3.375e-19}{J}\\
        E &= \frac{\num{3.375e-19}}{\num{1.602e-19}} = \qty{2.11}{eV} \quad\text{(J to eV)}}\\[3pt]
        This is less than $\phi = \qty{2.28}{eV}$.
  \item \eqlines{E &= \frac{hc}{\lambda} = \frac{\num{6.626e-34} \times \num{3.00e8}}{\num{2.00e-7}}
        = \qty{9.939e-19}{J} \quad\text{(photon energy)}\\
        E &= \frac{\num{9.939e-19}}{\num{1.602e-19}} = \qty{6.204}{eV} \quad\text{(J to eV)}\\
        K_\mathrm{max} &= hf - \phi = 6.204 - 4.70 = \qty{1.504}{eV}\\
        K_\mathrm{max} &= 1.504 \times \num{1.602e-19} = \qty{2.41e-19}{J} \quad\text{(eV to J)}\\
        V_0 &= \qty{1.50}{V} \quad\text{(the same number as $K_\mathrm{max}$ in eV)}\\
        v_\mathrm{max} &= \sqrt{\frac{2K_\mathrm{max}}{m_\mathrm{e}}}
        = \sqrt{\frac{2 \times \num{2.41e-19}}{\num{9.109e-31}}} = \qty{7.27e5}{m.s^{-1}}}\\[3pt]
        Answers: \final{$K_\mathrm{max} = \qty{1.50}{eV} = \qty{2.41e-19}{J}$},
        \final{$V_0 = \qty{1.50}{V}$}, \final{$v_\mathrm{max} = \qty{7.27e5}{m.s^{-1}}$}.
  \item (a) \eqlines{E &= \frac{hc}{\lambda} = \frac{\num{6.626e-34} \times \num{3.00e8}}{\num{4.05e-7}}
        = \qty{4.908e-19}{J} \quad\text{(one photon)}\\
        \frac{N}{t} &= \frac{P}{E} = \frac{\num{2.0e-3}}{\num{4.908e-19}} = \qty{4.075e15}{s^{-1}}
        \quad\text{(photons each second)}}\\[3pt]
        Answer: \final{\num{4.1e15} photons each second}.\\
        (b) In eV, one photon has
        \eqlines{E &= \frac{\num{4.908e-19}}{\num{1.602e-19}} = \qty{3.06}{eV} \quad\text{(J to eV)}}\\[3pt]
        This is more than $\phi = \qty{2.10}{eV}$, so electrons are released.
        Twice the power gives twice as many photons each second, so twice as many
        photoelectrons each second. Each photon still has the same energy, so
        $K_\mathrm{max} = \qty{3.06}{eV} - \qty{2.10}{eV} = \qty{0.96}{eV}$ does not
        change. On the current--voltage graph the saturation current \textbf{doubles},
        but the curve still meets the voltage axis at the same $-V_0$
        ($V_0 = \qty{0.96}{V}$).
  \item Q has the larger work function: its threshold frequency is bigger, and
        $\phi = hf_0$.\\[2pt]
        For P: \eqlines{\phi &= hf_0 = \num{6.626e-34} \times \num{6.0e14} = \qty{3.976e-19}{J}\\
        \phi &= \frac{\num{3.976e-19}}{\num{1.602e-19}} = \qty{2.48}{eV} \quad\text{(J to eV)}}\\[3pt]
        Answer: \final{$\phi_\mathrm{P} = \qty{2.5}{eV}$}.
  \item \textbf{Similarity} (1 mark): both said energy is given out or taken in only
        in packets (quanta) of $E = hf$. \textbf{Differences} (1 mark each):
        (1) Planck quantised the oscillators in a black body's walls; Einstein
        quantised light itself (photons). (2) Planck saw quanta as a maths trick;
        Einstein said photons are real. (Also accepted: they explained different
        experiments.)
\end{chapteranswers}
